\chapter{Green's Function Method for Cylindrically Symmetric Flux Tubes}
\label{ch:greensfunc}
\acresetall

\begin{summary}
	The effective action for arbitrary, non-homogeneous magnetic backgrounds can be expressed in terms 
	of Green's functions. In this chapter, I derive the differential equation 
	for these Green's functions for magnetic flux tubes of arbitrary profiles. This provides a 
	new tool for thinking about the flux tube effective action and suggests a new numerical 
	method that might be applied to the problem. 
	The Green's functions expressions are also useful for examining 
	field configurations with stationary effective actions. The specific case of step-function 
	flux tubes is discussed in detail in the Green's function picture.
\end{summary}

\section{Introduction}
In chapter \ref{ch:fluxtubes}, we discussed the astrophysical motivations for examining 
the effective actions of magnetic flux tubes in \ac{QED}. However, computing the fermion 
determinant for non-homogeneous magnetic fields can be quite involved. For computing magnetic 
flux tube effective actions, there are only a few methods available in the literature, some 
of which do not generalize easily to different flux tube profile shapes. These methods are 
discussed in section \ref{sec:EAisoflux}. Despite the calculational difficulties, 
the quantum correction to the classical energies of flux tube configurations are potentially 
significant, and could impact our conclusions about their stability or other properties. 
It is therefore interesting to explore other methods 
that might provide new tools for performing the calculation, and new insights into the problem.

In this chapter, I discuss some progress toward such a new method for the magnetic flux tube 
problem based on expressing the effective action in terms of integrals and sums over 
Green's functions. Since the Green's functions can be determined numerically, this 
method suggests a possible new numerical technique for computing effective actions. 
A Green's function technique has been used to exactly derive the effective action for a family of 
non-homogeneous magnetic fields in 2+1 dimensions~\cite{PhysRevD.52.R3163} and 3+1 dimensions
\cite{1998PhLB..419..322D}. The family of fields is given by 

\be
	\vec{B}(\vec{x}) = B_0 {\rm sech}{^2\left(\frac{x}{\lambda}\right)}\hat{z}.
\ee
For fields in this form, the effective action can be expressed as a finite integral. This 
result is particularly valuable because there are so few field configurations which are 
exactly solvable for the effective action. So, it provides a good tool for testing our 
understanding of the derivative expansion and other approximate expressions of the effective 
action. Presently, we will use the approach used to discover these exact solutions 
to explore the effective actions of 
flux tube configurations. 

\section{Effective Action in Arbitrary Fields}
\label{sec:arbitraryfields}

% The effective-action in the presence of a magnetic field is
% \be 
% i \int d^3x \mathcal{L} = {\rm ln ~Det} \{ D_\mu D^\mu + m^2 + \frac{1}{2}e\sigma^{\mu \nu}F^0_{\mu \nu} \}
% \ee 
% 
% where
% 
% \be 
% D_\mu = \partial_\mu + i e A_\mu
% \ee 
% and
% \be 
% \sigma^{\mu \nu} = \frac{i}{2}[\gamma^\mu,\gamma^\nu]
% \ee 
% 
 We begin with the effective action, equation (\ref{eqn:generalEA}), 
 in a static, but non-homogeneous magnetic field:
 
 \be
 \Gamma[A^0_\mu] = \Gamma_0 - \frac{i \hbar}{2} {\rm Tr ~ln} \left[\frac{(p_\mu+eA^0_\mu)^2
	+\frac{1}{2}e\sigma^{\mu\nu}F^0_{\mu\nu} + m^2}{p^2 +m^2}\right].
\ee
Since $A_\mu^0$ is time-independent and $A_0^0$ vanishes for static magnetic 
fields, we may express the canonical four-momentum in terms of its angular 
frequency eigenbasis and the canonical three-momentum,
\ba
	(p_\mu+eA^0_\mu)^2 &=& p^2 +e(p^\mu A^0_\mu + A^0_\mu p^\mu) + e (A^0)^2 \\
	& = & -\omega^2 -\vec{\nabla}^2 - ie[(\vec{\nabla}\cdot\vec{A}^0) + 
	(\vec{A}^0\cdot \vec{\nabla})] -e^2 (A^0)^2 \\
	& \equiv & -\omega^2 -\vec{\Pi}^2.
\ea
Then,

\be 
\label{eqn:EffAct}
\Gamma[A^0_\mu] = \Gamma_0 - \frac{i \hbar}{2} {\rm Tr ~ln} \left[\frac{-\omega^2 - \vec{\Pi}^2
	+\frac{1}{2}e\sigma^{\mu\nu}F^0_{\mu\nu} + m^2}{p^2 +m^2}\right] 
\ee
where
 \be 
	\sigma^{\mu \nu} = \frac{i}{2}[\gamma^\mu,\gamma^\nu]
 \ee 

and
\be 
	\Gamma_0 = \int d^4x\left(-\frac{1}{4}F^0_{\mu \nu}F^{0~\mu \nu}\right)
\ee 
is the classical action.

In order to express the effective action in terms of Green's functions, 
we make use of the fact that the trace involves an integration over 
frequency, $\omega$, and integration-by-parts in this integral 
produces the inverse of the operator:
\ba
	\lefteqn{{\rm Tr~ln}\left[-\omega^2-\vec{\Pi}^2+
	\frac{1}{2}e\sigma^{\mu \nu}F^0_{\mu \nu}+m^2\right]} \nonumber \\
	& = & \int_{-\infty}^\infty d\omega {\rm Tr}_3 ~{\rm ln}
	\left[-\omega^2-\vec{\Pi}^2+\frac{1}{2}e\sigma^{\mu \nu}F^0_{\mu \nu}+m^2\right] \\
	&=& 2\int_{-\infty}^\infty \omega^2 d\omega {\rm Tr}_3 
	\left[-\omega^2-\vec{\Pi}^2+\frac{1}{2}e\sigma^{\mu \nu}F^0_{\mu \nu}+m^2\right]^{-1} \\
	& = & \int_{-\infty}^\infty d\omega \int d^3\vec{x} d^3\vec{x}' 
	G_3(\vec{x},\vec{x}')\delta^{(3)}(\vec{x}-\vec{x}').
\ea

So, we may express (\ref{eqn:EffAct}) in terms of time-independent 
Green's functions to evaluate the trace:

\be
	\Gamma[A^0_\mu] =\Gamma_0 - i\hbar \int^{\infty}_{-\infty}\omega^2d\omega\int d^3\vec{x} 
		~\left(G_3(\vec{x},\vec{x})-G_3^0(\vec{x},\vec{x})\right),
\ee
where $G_3(\vec{x},\vec{x}')$ is defined by the Schr\"{o}dinger-like equation

\be 
\label{eqn:Greens}
\left[ -\vec{\Pi}^2 + \frac{1}{2}e\sigma^{\mu \nu}F^0_{\mu \nu} + m^2 
	-\omega^2\right]G_3(\vec{x},\vec{x}') 
	= \delta^{(3)}(\vec{x}-\vec{x}').
\ee 
$G_3^0(\vec{x},\vec{x}')$ is the Green's function in the absence of external fields (i.e. $A_\mu=0$),

\be 
\label{eqn:Green0}
\left[ -\vec{p}^2 + m^2 
	-\omega^2\right]G^0_3(\vec{x},\vec{x}') 
	= \delta^{(3)}(\vec{x}-\vec{x}').
\ee 

\section{Cylindrically Symmetric B-fields}

We now specialize our analysis to the case of a cylindrically symmetric magnetic field, 
pointing in the \bhattext{z}-direction in cylindrical coordinates, the symmetry relevant 
to flux tube configurations.
Making a gauge choice that $A_0 = A_\rho = A_z = 0$, 
the magnetic field is $\vec{B} = B(\rho)$\bhattext{z} with

\be 
\label{eqn:Bgf}
	B(\rho) = \frac{A_\phi(\rho)}{\rho} + \frac{dA_\phi(\rho)}{d\rho}.
\ee 


Using the cylindrical symmetry, the Green's function is a function of the radial 
coordinate and can be expressed as a sum over the magnetic quantum number, $m_l$, 
the momentum along the axis of symmetry, $k_z$, the frequency, $\omega$, and 
the spin projection eigenvalues, $\sigma^3$:

\be
	G_3(\vec{x},\vec{x}') = \sum_{m_l,k_z,\omega,\sigma^3} G_{m_l,k_z,\omega,\sigma^3}(\rho,\rho').
\ee
From equation (\ref{eqn:Greens}), we can write the equation governing each term of the 
Green's function,

\be 
\label{eqn:cylGreen}
\left[ -\frac{d^2}{d\rho^2}  - \frac{1}{\rho}\frac{d}{d\rho}+ \frac{1}{\rho^2}+V_{m_l}(\rho) +k_z^2 +m^2-\omega^2
	 +\frac{1}{\rho^2}\right]\cylGreen 
	= \delta(\rho-\rho')
\ee 
with the magnetic potential given by

\be 
V_{m_l}(\rho) = \frac{1}{2}e\sigma^{\mu \nu}F^0_{\mu \nu} +\frac{(m_l^2-1)}{\rho^2} 
	-\frac{2m_l e}{\rho}\Aphi +e^2 (\Aphi)^2.
\ee 

The spin-field coupling is proportional to the spin eigenvalue and the magnetic field,
$\sigma^{\mu \nu}F_{\mu\nu} = 2 \sigma^{1 2}B_z = 
	2 \sigma ^3\left(\frac{\Aphi}{\rho} +\frac{d\Aphi}{d\rho}\right)$.  So, the 
potential term is given by


%  in a coordinate system with metric 
% signature ${\rm diag}(-1, 1, \rho^2, 1)$.  Letting Roman characters label the local Lorentz  
% coordinate system and Greek letters label the Cylindrical coordinates, our tetrads, defined by the equation 
% $g^{\mu \nu} = e^\mu_a e^\nu_b \eta^{a b}$ are given by
% 
% \ba 
% e^\phi_2  =\rho \nonumber \\
% e^t_0 = e^\rho_1 =  e^z_3 = 1  \nonumber \\
% e^\mu_a = 0 ~~{\rm otherwise}
% \ea 
% We may now write our gamma matrices in cylindrical coordinates using the relation
% 
% \be 
% \gamma^\mu = e^\mu_a \gamma^a
% \ee 
% so that
% 
% \ba 
% \label{eqn:cylgammas}
% \gamma^\rho = \gamma^1 \nonumber \\
% \gamma^\phi = \rho \gamma^2 \nonumber \\
% \gamma^z = \gamma^3
% \ea where the numerical superscripts refer to standard Minkowski space $\gamma$ matrices.
% 
% Using the results (\ref{eqn:cylgammas}), we have
% 
% \ba 
% \label{eqn:sigmaF}
% \sigma^{\mu \nu}F^0_{\mu \nu} &=& \sigma^{\rho \phi} F^0_{\rho \phi} + 
% 		\sigma^{\phi \rho } F^0_{\phi \rho} \nonumber \\
% 	& = & 2 \rho \sigma^{3} \pderiv{\Aphi}{\rho}.
% \ea  with $\sigma^3$ denoting the spin projection eigenvalue, $+/- 1$ for spin up/down.
% 
% Putting equation (\ref{eqn:sigmaF}) into equation (\ref{eqn:GreensCyl}) gives

\be 
\label{eqn:effPot}
V_{m_l}(\rho) = e\sigma^{3} \left(\frac{\Aphi}{\rho}  + 
	\frac{d\Aphi}{d\rho}\right) + \frac{(m_l^2-1)}{\rho^2}+
	e^2 (\Aphi)^2 - \frac{2e m_l}{\rho} \Aphi.
\ee 
$\sigma^3$ is the spin projection eigenvalue $+1/-1$ for spin up/down.

\section{Solution Strategy}



We will attempt to find the effective action by computing the Green's function, 
(\ref{eqn:cylGreen}).  This can be done by
matching the independent solutions of the homogeneous equation,

\be
\left[ -\frac{d^2}{d\rho^2}  - \frac{1}{\rho}\frac{d}{d\rho}+ \frac{1}{\rho^2}+V_{m_l}(\rho) +k_z^2 +m^2-\omega^2
	 +\frac{1}{\rho^2}\right]u_{m_l,k_z,\omega,\sigma^3}(\rho) 
	= 0.
\ee
For a given model of the magnetic field (a choice of $\Aphi$), 
the homogeneous solutions $u_0(\rho)$ and $u_\infty(\rho)$ are regular at $\rho =0$
and $\rho=\infty$ respectively.  The Wronskian of these solutions is defined as

\be 
W(\rho) = u'_0(\rho)u_\infty(\rho) - u_0(\rho)u_\infty'(\rho).
\ee 
Then, the Green's function which solves (\ref{eqn:cylGreen}) is simply given by

\be 
G_{\omega,k_z, m_l} (\rho, \rho') = \theta(\rho' - \rho) \frac{\rho u_0(\rho)u_\infty(\rho')}{W_0} + 
	\theta(\rho-\rho')\frac{\rho u_0(\rho')u_\infty(\rho)}{W_0}
\ee 
where
\be 
W_0 = \rho W(\rho)
\ee 
is constant by Abel's theorem.

\begin{sloppypar}
We can exactly solve for the Green's function in the 
absence of fields, $G_{\subscripts}^0(\rho,\rho)$,
 by noticing that the 
differential equation is Bessel's equation in the dimensionless 
coordinate $x=\sqrt{\omega^2-k_z^2 -m^2}\rho$.
Then, the homogeneous solution that is regular at $\rho=0$, $u_0(\rho)$ is identified 
with $J_{m_l}(\sqrt{\omega^2-k_z^2-m^2}\rho)$ and the homogeneous solution 
that is regular at $\rho=\infty$ is identified with $Y_{m_l}(\sqrt{\omega^2-k_z^2-m^2}\rho)$, 
where $J_n(x)$ and $Y_n(x)$ are Bessel functions of the first and second kind,
respectively.
So, we have

\be 
G_{\subscripts}^0(\rho,\rho) = \frac{\rho}{W^0_0}J_{m_l}(\sqrt{\omega^2-k_z^2-m^2}\rho)
	Y_{m_l}(\sqrt{\omega^2-k_z^2 -m^2}\rho).
\ee 
We can also compute
$W_0^0$ exactly for the field-free case using the Bessel function identity

\be 
\label{eqn:BessID}
J'_n(x)Y_n(x) - J_n(x)Y'_n(x)=-\frac{2}{\pi x}.
\ee 
Then,
\be 
\label{eqn:BkgGreens}
G_{\subscripts}^0(\rho,\rho) = -\frac{\pi \rho}{2}J_{m_l}(\sqrt{\omega^2-k_z^2 -m^2}\rho)
	Y_{m_l}(\sqrt{\omega^2-k_z^2 -m^2}\rho)
\ee 
\end{sloppypar}
We can now write the (unrenormalized) effective action as

\ba 
	\Gamma &=& \Gamma_0 -i\hbar \sum_{\rm \sigma^3= \{\pm1\}}\sum_{m_l=-\infty}^\infty  
	 \int_{-\infty}^\infty \omega^2 d\omega dk_z
	\int_0^{\infty} d\rho \rho (G_{\subscripts}(\rho,\rho) \nonumber \\ 
	& &~~~~~~- G_{\subscripts}^0(\rho,\rho)) \\
\label{eqn:soln}
	&=& \Gamma_0 -i\hbar \sum_{\rm \sigma^3= \{\pm1\}}\sum_{m_l=-\infty}^\infty  
	\int_{-\infty}^\infty \omega^2 d\omega dk_z
	\int_0^{\infty} d\rho \rho^2 \biggr(\biggr[\frac{ 
	u_0(\rho)u_\infty(\rho)}{W_0}\biggr]_{\subscripts} \nonumber \\ 
	& &~~~~~~~+\frac{\pi}{2}J_{m_l}(\sqrt{\omega^2-k_z^2 -m^2}\rho)
	Y_{m_l}(\sqrt{\omega^2-k_z^2 -m^2}\rho)\biggr) \nonumber \\
	&=& \Gamma_0 +\hbar\pi \sum_{\rm \sigma^3= \{\pm1\}}\sum_{m_l=-\infty}^\infty  
	\int_{0}^\infty \chi^3 d\chi
	\int_0^{\infty} d\rho \rho^2 \biggr(\biggr[\frac{ 
	u_0(\rho)u_\infty(\rho)}{W_0}\biggr]_{\csubscripts} \nonumber \\ 
	& &~~~~~~~+\frac{\pi}{2}J_{m_l}(\sqrt{\chi^2 -m^2}\rho)
	Y_{m_l}(\sqrt{\chi^2 -m^2}\rho)\biggr).
\ea 
The final equality is arrived at by analytically 
continuing $\omega\to i\omega$ and observing that 
the value of the integral along a semi-circle which 
closes the contour is unchanged by the rotation 
in the complex plane.  
By Abel's identity, $W_0$ has no dependence on $\rho$.  
It can come outside of the integral.  In the above equations, we have introduced a sum 
over the eigenvalues of the operator $\sigma^{1 2} = \frac{i}{2}(\gamma^1 \gamma^2 - \gamma^2 \gamma^1)$.
So, to account for two pairs of degenerate eigenvalues, 
we introduce an overall factor of $2$ and sum over $\sigma^3=\{+1,-1\}$.

The $u(\rho)_{\csubscripts}$'s are solutions to the differential equation

\be 
\label{eqn:chieqn}
\left(-\frac{d^2}{d\rho^2} - \frac{1}{\rho}\frac{d}{d\rho} + 
	V_{m_l}(\rho) - \chi^2 + m^2 +\frac{1}{\rho^2}\right)u_{[0,\infty]}(\rho)|_{\chi,m_l,\sigma^3} = 0
\ee 
where $[u_0(\rho)]_{\csubscripts}$ and $[u_\infty(\rho)]_{\csubscripts}$ are independent solutions which 
are regular at $\rho = 0$ and $\rho = \infty$, respectively.

% We may simplify the following numerics by absorbing the $\chi^2-m^2$ 
% term into an order-0 Bessel function so that
% 
% \be 
% [u_0(\rho)]_{\csubscripts} = J_0(\sqrt{\chi^2-m^2}\rho)[\psi_0(\rho)]_{m_l,\sigma^3}
% \ee 
% \be 
% [u_\infty(\rho)]_{\csubscripts} = Y_0(\sqrt{\chi^2-m^2}\rho)[\psi_\infty(\rho)]_{m_l,\sigma^3}
% \ee 
% where the $\psi$ functions satisfy
% \be 
% \left[\frac{d^2}{d\rho^2} +\left(\frac{1}{\rho}-
% 	2\sqrt{\chi^2-m^2}\frac{J_1(\sqrt{\chi^2-m^2}\rho)}{J_0(\sqrt{\chi^2-m^2}\rho)}\right)\frac{d}{d\rho}
% 	-V_{m_l}(\rho) -\frac{1}{\rho^2}\right]
% 	[\psi_0(\rho)]_{m_l,\sigma^3}=0
% \ee 
% \be 
% \left[\frac{d^2}{d\rho^2} +\left(\frac{1}{\rho}-
% 	2\sqrt{\chi^2-m^2}\frac{Y_1(\sqrt{\chi^2-m^2}\rho)}{Y_0(\sqrt{\chi^2-m^2}\rho)}\right)\frac{d}{d\rho}
% 	-V_{m_l}(\rho) - \frac{1}{\rho^2}\right]
% 	[\psi_\infty(\rho)]_{m_l,\sigma^3}=0
% \ee 
% 
% The constant $W_0$ can be expressed in terms of the $\psi$ functions
% \be 
% W_0=-\frac{2}{\pi}[\psi_0(\rho)\psi_\infty(\rho)]_{m_l,\sigma^3} + \rho W_\psi(\rho)
% 	J_0(\sqrt{\chi^2-m^2}\rho)Y_0(\sqrt{\chi^2-m^2}\rho)
% \ee 
% where $W_\psi$ has the form of a Wronskian, but it should be noted that $\psi_0$ and $\psi_\infty$ 
% are solutions to different differential equations.
% \be 
% W_\psi(\rho) = \psi'_0(\rho)\psi_\infty(\rho) - \psi_0(\rho)\psi'_\infty(\rho).
% \ee 

% The effective action is an extensive quantity.  Therefore, we divide by the (infinite) spacetime volume to 
% find the effective vacuum potential.
% 
% \be 
% 	V_{eff}(\Aphi) = \frac{\Gamma}{\pi L_{\rho}^2 \int_{-\infty}^{\infty} d\omega dk_z}
% \ee 
% 

Equations (\ref{eqn:soln}) and (\ref{eqn:chieqn}) are the most general form for the 
effective action in this method. To use these equations for computations, we must know the 
functions $[u_0(\rho)]_{\csubscripts}$ and $[u_\infty(\rho)]_{\csubscripts}$. Therefore, 
we can proceed either by choosing a form of the potential, $V_{m_l}(\rho)$, and 
solving the differential equation analytically, or by computing the required functions 
numerically for an arbitrary potential. In section \ref{sec:stepfunc}, we will 
make some analytic progress using a step-function flux tube profile. For either 
technique, though, we must establish the boundary conditions on these solutions, 
so we discuss those next.



\section{Boundary Conditions}
We would like to choose suitable boundary conditions to numerically 
compute $[u_0(\rho)]_{\csubscripts}$ and $[u_\infty(\rho)]_{\csubscripts}$.  We begin by 
considering $\Aphi$ in the form

\be 
\Aphi = \frac{F}{2\pi \rho}f_\lambda(\rho)
\ee 
so that
\be 
B_z(\rho)=\frac{F}{2\pi\rho}\frac{df_\lambda(\rho)}{d\rho}.
\ee 
The subscript $\lambda$ refers to a length scale that characterizes the size of the flux 
tube.  In practice, it is inconvenient to discuss the boundary conditions at infinity. So, 
we will introduce a scale, $L_\rho \gg \lambda$, which is finite but 
sufficiently far away from the flux tube. 
The total magnetic flux is
\be 
\Phi=F(f_\lambda(L_\rho)-f_\lambda(0)).
\ee 
It is convenient to express the flux in units of $\frac{2 \pi}{e}$ and define 
a dimensionless quantity

\be 
\mathcal{F}=\frac{e}{2 \pi} F
\ee 
Under these definitions, our operator is

\be 
-\frac{d^2}{d\rho^2} - \frac{1}{\rho}\frac{d}{d\rho} + 
	\frac{\sigma^3\mathcal{F}}{\rho}f'_\lambda(\rho) +
	 \frac{1}{\rho^2}\left(m_l -\mathcal{F}f_\lambda(\rho)\right)^2
	-\left(\chi^2-m^2\right).
\ee 

If we require that $B_z(\rho)$ remains finite at $\rho \to 0$ and that
 $\frac{dB_z(\rho)}{d\rho}\biggr|_{\rho=0}=0$, then we can choose without 
loss of generality that, for small $\rho$, $f_\lambda(\rho) = C_1\rho^2 + C_2\rho^4 +...$ 
and $f_\lambda(L_\rho) = 1$. Then, the total flux, $\Phi$, is simply $F$, and the 
field profile near the center is a quadratic with coefficient given by

\be 
C_1=\frac{eB_z(0)}{2 \mathcal{F}}\equiv\frac{e B_0}{2\mathcal{F}}.
\ee 

With the above observations, we notice that equation (\ref{eqn:chieqn}) takes the form of 
Bessel's equation in the asymptotic limits of small and large $\rho$.  
We use this fact to 
inform the following boundary conditions:


\ba 
\label{eqn:BCa}
 u_0(0)|_{\csubscripts} &=& J_{m_l}(0) \\
 u'_0(0)|_{\csubscripts} &=& \sqrt{\chi^2-m^2-eB_0(\sigma^2-m_l)} J'_{m_l}(0)\nonumber \\
	&=&\frac{\sqrt{\chi^2-m^2-eB_0(\sigma^3-m_l)}}{2} \nonumber \\
 	& & \times(J_{m_l-1}(0)-J_{m_l+1}(0))\\
 u_\infty(L_{\rho})|_{\csubscripts} &=& Y_n(\sqrt{\chi^2-m^2}L_\rho) \\
\label{eqn:BCb}
 u'_\infty(L_{\rho})|_{\csubscripts} &=& \sqrt{\chi^2-m^2-eB_0(\sigma^2-m_l)} Y'_{m_l}(0)\nonumber \\
	&=&\frac{\sqrt{\chi^2-m^2}}{2}
  	(Y_{n-1}(\sqrt{\chi^2-m^2}L_\rho) \nonumber \\
 	& & -Y_{n+1}(\sqrt{\chi^2-m^2}L_\rho)) \\
 n&\equiv&m_l-\mathcal{F}.
\ea 

If we take $F=0$, these boundary conditions reproduce our result for the 
Green's function in the absence of any fields, equation (\ref{eqn:BkgGreens}).  

% In terms of the $\psi(\rho)$ functions, our boundary conditions are
% \ba 
% \label{eqn:BCa}
%  \psi_0(0)|_{m_l,\sigma^3} &=& \frac{J_{m_l}(0)}{J_0(0)} \\
%  \psi'_0(0)|_{m_l,\sigma^3} &=& \frac{\sqrt{\chi^2-m^2-eB_0(\sigma^3-m_l)}}{2}
%  	\frac{J_{m_l-1}(0)-J_{m_l+1}(0)}{J_0(0)}\\
%  \psi_\infty(L_{\rho})|_{m_l,\sigma^3} &=& 
% 	\frac{Y_n(\sqrt{\chi^2-m^2}L_\rho)}{Y_0(\sqrt{\chi^2-m^2})} \\
%  \psi'_\infty(L_{\rho})|_{m_l,\sigma^3} &=& \frac{\sqrt{\chi^2-m^2}}{2} \nonumber \\
%   	& & \frac{Y_{n-1}(\sqrt{\chi^2-m^2}L_\rho) 
%  	-Y_{n+1}(\sqrt{\chi^2-m^2}L_\rho)}{Y_0(\sqrt{\chi^2-m^2}L_\rho)} \\
% \label{eqn:BCb}
% \ea 



% We would like to choose suitable boundary conditions to numerically 
% compute $[u_0(\rho)]_{\subscripts}$ and $[u_\infty(\rho)]_{\subscripts}$.  
% Consider the limit of asymptotically large $\rho$.  As a consequence of (\ref{eqn:B}),
% for magnetic fields localized 
% around finite $\rho$, we must have $\Aphi \rightarrow {\rm const.}$ as $\rho \rightarrow \infty$.
% Thus, we recover Bessel's equation in this limit, allowing us to define boundary conditions 
% that guarantee the linear independence of our solutions.
% 
% \ba 
% u_0(\rho=L_{\rho}) &=& J_{m_l}(sL_{\rho}) \\
% u'_0(\rho=L_{\rho}) &=& J'_{m_l}(sL_{\rho}) \\
% u_\infty(\rho=L_{\rho}) &=& Y_{m_l}(
% 	sL_{\rho}) \\
% u'_\infty(\rho=L_{\rho}) &=& Y'_{m_l}(sL_{\rho}) \\
% s&\equiv&\sqrt{\omega^2 -k_z^2-m^2  
%  	+ (eA_\phi(L_\rho))^2}
% \ea 
% where $L_\rho$ is chosen relative to a length scale, $\lambda$, characterizing 
% the flux tube, $\lambda \ll L_\rho$.  Here, we will take $A_\phi(L_\rho)=0$, 
% which must be true if we are to have finite total flux.
% 
% Making use of (\ref{eqn:BessID}) and this particular choice of boundary conditions 
% we see that 
% 
% \be 
% 	W_0^{-1}=-\frac{\pi}{2}\sqrt{\omega^2 -k_z^2-m^2}
% \ee 
% 
% 
% Now,
% \ba 
% \label{eqn:soln}
% 	\Gamma&=& \Gamma_0 +\frac{i \hbar \pi}{2} 
% 	\sum_{\rm \sigma^3= \{\pm1\}}\sum_{m_l=-\infty}^\infty  
% 	\int_{-\infty}^\infty \omega^2 d\omega dk_z
% 	\int_0^{L_\rho} d\rho \rho^2 \nonumber \\
% 	& & \times \sqrt{\omega^2 -k_z^2 -m^2}\biggr(\biggr[ 
% 	u_0(\rho)u_\infty(\rho)\biggr]_{\subscripts} \nonumber \\ 
% 	& &- 
% 	J_{m_l}(\sqrt{\omega^2-k_z^2 -m^2}\rho)
% 	Y_{m_l}(\sqrt{\omega^2-k_z^2 -m^2}\rho)\biggr)
% \ea 
% We may analytically continue $k_z\rightarrow ik_z$ as usual\footnote{The value of the 
% integral along a semi-circle that closes the contour is non-zero but is the same 
% before and after the rotation, so we may neglect it.} and 
% write the resulting integral in polar coordinates in the Wick rotated $\omega$-$k_z$ plane.  
% Integrating over the angular part leaves
% \ba 
% 	\Gamma &=& \Gamma_0 -\frac{\hbar \pi^2}{2} 
% 	\sum_{\rm \sigma^3= \{\pm1\}}\sum_{m_l=-\infty}^\infty  
% 	\int_{0}^\infty  d\chi \chi^3
% 	\int_0^{L_\rho} d\rho \rho^2 \sqrt{\chi^2 -m^2 }\nonumber \\
% 	& & \times \biggr(\biggr[ 
% 	\hat{u}_0(\rho)\hat{u}_\infty(\rho)\biggr]_{\chi,m_l,\sigma^3} - 
% 	J_{m_l}(\sqrt{\chi^2 -m^2}\rho)
% 	Y_{m_l}(\sqrt{\chi^2 -m^2}\rho)\biggr).
% \ea 
% The $\hat{u}(\rho)$'s are solutions to the differential equation
% 
% \be 
% \left(-\frac{d^2}{d\rho^2} - \frac{1}{\rho}\frac{d}{d\rho} + 
% 	V_{m_l}(\rho) - \chi^2 +\frac{1}{\rho^2}\right)\hat{u}_{[0,\infty]}(\rho)|_{\chi,m_l,\sigma^2} = 0
% \ee 

From equation (\ref{eqn:soln}), we can compute the effective action given the homogeneous solutions 
to the Green's function, $u_0(\rho)$ and $u_\infty(\rho)$.
Equations (\ref{eqn:cylGreen}) and (\ref{eqn:effPot}) define an \ac{ODE} that can be numerically 
integrated to find the homogeneous solutions given a cylindrically symmetric vector potential $A_\phi(\rho)$
and suitable boundary conditions, equations (\ref{eqn:BCa}) to (\ref{eqn:BCb}).

\section{Configurations with Stationary Action}
\label{sec:stationaryEA}

The classical motion of a system is determined by the principle of least action: 
classical systems follow paths for which the variation of the action is zero. 
\be 
\delta \Gamma = 0
\ee 
The quantum 
mechanical explanation for this is well-known. Consider the partition functional 
of \ac{QED}, equation (\ref{eqn:qedpartition}),

\ba
\lefteqn{Z[J^\mu, \bar{\eta}, \eta]}\nonumber \\
  & =&\int \mathcal{D}A_\mu \mathcal{D}\psi \mathcal{D} \bar{\psi}
  \exp\left( \frac{i}{\hbar}\int d^4x(\mathcal{L} + J^\mu A_\mu + \bar{\eta}\psi
  +\bar{\psi}\eta)\right).
  \label{eqn:qedpartition2}
\ea
In a classical system, the action in units of $\hbar$ is very large. This means 
that the functional integral is highly oscillatory and for most field configurations 
gives a small contribution. Therefore, the main contributions come from critical 
points where the action is stationary relative to the fields. These critical points 
are simply the solutions to the classical field equations. This observation is the 
basis of the stationary phase method in quantum mechanics, which replaces the functional 
integrals with a sum of integrals over neighbourhoods of stationary action. 

Since the 
effective action can be thought of as a quantum-correction of the classical action, 
the principle of least effective action would give us the quantum-corrected 
equations of motion. The Green's function method allows us to express the 
effective action in terms of functionals of an arbitrary magnetic field profile.
So, we take this opportunity to explore the stationary points of the 
effective action with respect to the flux tube profile shape.


In our geometry, the classical action is

\be 
\Gamma_0=-\int d^4x \frac{B^2}{2} = -\pi \int_{-\infty}^\infty dt dz \int_0^{L_\rho} \rho d\rho
	B_z(\rho)^2.
\ee We would like to minimize the effective action with respect to fluctuations 
of the function $\flr$ and its first derivative $f'_\lambda(\rho)$.  
Taking the entire action, 

\be 
\delta \Gamma = \pderiv{\Gamma}{\flr} \delta \flr + \pderiv{\Gamma}{(f'_\lambda(\rho))}\delta f'_\lambda(\rho).
\ee 

For now, we will neglect the vacuum renormalization terms and reintroduce them at the end of the analysis:

\ba 
\label{eqn:varActa}
\delta \Gamma &=& \biggr[\biggr\{-2\pi \int_{-\infty}^\infty dt dz B_z(\rho) + 
	\hbar \pi \sum_{m_l,\sigma^3} \int_{0}^\infty \chi^3 d\chi W_0^{-1}\nonumber \\
	& &\left(
	\pderiv{u_0(\rho)}{f'_\lambda(\rho)}u_\infty(\rho) +u_0(\rho)\pderiv{u_\infty(\rho)}{f'_\lambda(\rho)}
	\right)\biggr\}\delta \flr \biggr]_{\rho=0}^{L_\rho} \nonumber \\
	& &+\int_0^{L\rho}d\rho\biggr\{2\pi\int_{-\infty}^\infty dtdz\left(\frac{dB_z(\rho)}{d\rho}\right) 
	-\hbar \pi\sum_{m_l,\sigma^3} \int_{0}^\infty \chi^3d\chi 
	 \rho^2 W_0^{-1}\nonumber \\
	& & \times \biggr[u_\infty(\rho)\left(\pderiv{u_0(\rho)}{\flr} - \pderiv{u'_0(\rho)}
	{f'_\lambda(\rho)}\right) + u_0(\rho)\left(\pderiv{u_\infty(\rho)}{\flr} - \pderiv{u'_\infty(\rho)}
	{f'_\lambda(\rho)}\right) \nonumber \\
	& & -\pderiv{u_0(\rho)}{f'_\lambda(\rho)}u'_\infty(\rho) -
	 \pderiv{u_\infty(\rho)}{f'_\lambda(\rho)}u'_0(\rho)
	\biggr]_{\chi, m_l, \sigma^3}\biggr\}\delta\flr = 0.
\ea 

To compute the functional derivatives of the $u(\rho)$'s, we make use of the chain rule.  Here it is 
appropriate since functionals can be defined in terms of integrals over composed functions~\cite{greiner1996field}:

\be 
\frac{\delta u(\rho, [\flr])}{\delta \flr} = \pderiv{u(\rho)}{\rho}\frac{d\rho}{d \flr} = \frac{u'(\rho)}{f'_\lambda(\rho)}.
\ee 

The surface terms for the classical part vanish straightforwardly.  
We can impose conservation of flux by fixing the 
endpoints of the field fluctuations so that $\delta f_\lambda(0)=\delta f_\lambda(L_\rho) =0$.  
In this case, we can ignore the surface terms in (\ref{eqn:varActa}) to arrive at




\ba 
\label{eqn:varAct}
\int_0^{L\rho}d\rho\biggr\{2\pi\int_{-\infty}^\infty dtdz\left(\frac{d B_z(\rho)}{d\rho}\right) 
	-\frac{\hbar \pi}{2}\sum_{m_l,\sigma^3} \int_{0}^\infty \chi^3d\chi 
	 \rho^2  W_0^{-1}& &\nonumber \\
	 \times \biggr[u_\infty(\rho)\left(\frac{u'_0(\rho)}{f'_\lambda(\rho)} - \frac{u''_0(\rho)}
	{f''_\lambda(\rho)}\right) + u_0(\rho)\left(\frac{u'_\infty(\rho)}{f'_\lambda(\rho)} 
	- \frac{u''_\infty(\rho)}
	{f''_\lambda(\rho)}\right) & &\nonumber \\
	 -2\frac{u'_0(\rho)u'_\infty(\rho)}{f''_\lambda(\rho)}
	\biggr]_{\chi,m_l,\sigma^3}\delta \flr\biggr\} = 0. & &
\ea 
Inside the integral, there are no restrictions on the fluctuations $\delta \flr$
meaning that the rest of the integrand must vanish on its own:



\ba 
\label{eqn:stationary}
\frac{d B_z(\rho)}{d \rho}
	-\frac{\hbar}{2\Delta} \int_{0}^\infty \chi^3d\chi 
	 \rho^2 
	 \sum_{m_l,\sigma^3}W_0^{-1}\biggr[u_\infty(\rho)\biggr(\frac{u'_0(\rho)}{f'_\lambda(\rho)} 
	- \frac{u''_0(\rho)}{f''_\lambda(\rho)}\biggr) 
	& & \nonumber \\
	+ u_0(\rho)\left(\frac{u'_\infty(\rho)}{f'_\lambda(\rho)} 
	- \frac{u''_\infty(\rho)}
	{f''_\lambda(\rho)}\right) 
	 -2\frac{u'_0(\rho)u'_\infty(\rho)}{f''_\lambda(\rho)}
	\biggr]_{\chi,m_l,\sigma^3} =0,
\ea 
where $\Delta=\int_{-\infty}^{\infty}d\omega \int_{-\infty}^{\infty}dk_z$ is an infinite constant
arising because the field extends infinitely in time and along the $\hat{z}$-direction.

Finally, we subtract the zero flux 
equivalents of each of the one-loop terms to ensure that the action 
remains finite and vanishes with the magnetic field. Thus, we arrive at the 
 quantum-corrected equations of motion for a magnetic flux tube:

 \ba 
 \label{eqn:stationaryEffAct}
 &&\frac{d B_z(\rho)}{d \rho}
 	-\frac{\hbar}{2\Delta} \int_{0}^\infty \chi^3d\chi 
 	 \rho^2 
 	 \sum_{m_l,\sigma^3}\biggr[W_0^{-1}\biggr\{u_\infty(\rho)\biggr(\frac{u'_0(\rho)}{f'_\lambda(\rho)} 
 	- \frac{u''_0(\rho)}{f''_\lambda(\rho)}\biggr) 
 	 \nonumber \\
 	& &+ u_0(\rho)\left(\frac{u'_\infty(\rho)}{f'_\lambda(\rho)} 
 	- \frac{u''_\infty(\rho)}
 	{f''_\lambda(\rho)}\right) 
 	 -2\frac{u'_0(\rho)u'_\infty(\rho)}{f''_\lambda(\rho)}
 	\biggr\} \nonumber \\
 	& &+\frac{\pi}{2}\biggr\{\BESSO{Y}\biggr(\frac{\sqrt{\chi^2-m^2}\BESSO{J'}}{f'_\lambda(\rho)} \nonumber \\
 	& &- \frac{(\chi^2-m^2)\BESSO{J''}}{f''_\lambda(\rho)}\biggr) 
 	\nonumber \\
 	& & + \BESSO{J}\biggr(\frac{\sqrt{\chi^2-m^2}\BESSO{Y'}}{f'_\lambda(\rho)} \nonumber \\
 	& &- \frac{(\chi^2-m^2)\BESSO{Y''}}
 	{f''_\lambda(\rho)}\biggr)  \nonumber \\
 	& & -2\frac{(\chi^2-m^2)\BESSO{J'}\BESSO{Y'}}{f''_\lambda(\rho)}
 	\biggr\}
 \biggr]_{\chi,m_l,\sigma^3} =0.
 \ea 
 This expression is not guaranteed to be finite because it does not yet account for the field-strength 
 renormalization. It is not clear how this renormalization can be performed for the general case.

% In terms of the $\psi$ fields, this result becomes
% \ba 
% & & \frac{d B_z(\rho)}{d \rho}
% 	-\frac{\hbar}{2\Delta} \sum_{m_l,\sigma^3}\biggr[\int_{0}^\infty \chi^3 
% 	\frac{J_0(\sqrt{\chi^2-m^2}\rho)Y_0(\sqrt{\chi^2-m^2}\rho)}{W_0} d\chi \biggr]
% 	 \rho^2 \nonumber \\
% 	 & &\biggr[\psi_\infty(\rho)\biggr(\frac{\psi'_0(\rho)}{f'_\lambda(\rho)} 
% 	- \frac{\psi''_0(\rho)}{f''_\lambda(\rho)}\biggr) 
% 	+ \psi_0(\rho)\left(\frac{\psi'_\infty(\rho)}{f'_\lambda(\rho)} 
% 	- \frac{\psi''_\infty(\rho)}
% 	{f''_\lambda(\rho)}\right) \nonumber \\
% 	 & & -2\frac{\psi'_0(\rho)\psi'_\infty(\rho)}{f''_\lambda(\rho)}
% 	\biggr]_{m_l,\sigma^3} =0
% \ea 

Here we see the familiar classical result:  the field which minimizes the classical action is homogeneous 
($\frac{d B_z(\rho)}{d \rho} =0$ everywhere).  However, the second term arising from the one-loop 
QED effects suggests that some non-homogeneous field ($\frac{d B_z(\rho)}{d \rho} \ne 0$)
may be a stationary point of the action for a fixed amount of flux. This equation 
potentially provides a numerical method for computing the quantum-corrected equations 
of motion of a cylindrically symmetric magnetic field. However, exploring 
this equation in more depth is outside the scope of this dissertation.
 

%\section{Flux Tube Models}

%The macroscopic properties of superconductors are usually described by 
%the Landau-Ginzburg effective potential.  In this model, the gauge field 
%which describes the flux tube at large distances is 

%\be 
%\Aphi = \frac{1}{e}\left(\frac{1}{\rho} + \frac{C}{\lambda} K_1(\rho/\lambda)\right)
%\ee where $C$ is an arbitrary constant governing the 
%magnetic flux and $\lambda$ is the London 
%penetration length

%\be 
%\lambda = \sqrt{\frac{m}{4 \pi e^2 n}}.
%\ee 
%Here, $n$ is the number density of Cooper pairs in the superconducting medium.


%Another model magnetic field profile is the Gaussian flux tube. We have 

%\be 
%B_z(\rho) = \frac{F}{\pi \lambda^2} \exp{(-(\rho/\lambda)^2)}
%\ee 
%A gauge field which gives us this profile is

%\be 
%\Aphi = \frac{F}{2 \pi \rho}(1-\exp{(-\rho^2/\lambda^2)})
%\ee 
%The Gaussian is particularly appealing since the field, $\Aphi$, and its derivatives are finite 
%as $\rho \to 0$.

%If we wish to have a Step-Function profile, we have 

%\be 
%B_z(\rho) = \frac{F}{\pi \lambda^2} \theta(\lambda-\rho) =\frac{2 \mathcal{F}}{e\lambda^2}\theta(\lambda-\rho)
%\ee 
%A gauge field which gives us this profile is

%\be 
%\Aphi = \frac{F}{2\pi} \left(\frac{\rho}{\lambda^2}\theta(\lambda-\rho) +
%	\frac{1}{\rho}\theta(\rho-\lambda)\right)
%\ee 

%For a flux tube with a solitonic profile, we have 

%\be 
%B_z(\rho) = \frac{F}{2 \pi cosh^2(\rho/\lambda)}
%\ee 
%A gauge field which gives us this profile is

%\be 
%\Aphi = \frac{F}{2\pi} \left( \lambda \tanh{(\rho/\lambda)}.-\frac{\lambda^2}{\rho}\ln(\cosh{(\rho/\lambda)})\right)
%\ee 

%A zero flux configuration explicitly accounts for the back-flux and has the desirable property
%that $\Aphi$ falls faster than $1/\rho$ for large $\rho$.  However, we must replace the 
%boundary condition at $L_\rho$ with $f_\lambda(L_\rho) = 0$.  An example zero-flux field is
%\be 
%\Aphi = \frac{F}{2 \pi} \frac{\rho^3}{e^{(\rho/\lambda)^2}-1}.
%\ee 

\section{Step-Function Flux Tube}
\label{sec:stepfunc}
As an example of how the Green's function method is applied to specific magnetic field 
profiles, $f_\lambda(\rho)$, we consider a step-function flux tube. The \ac{QED} effective 
action in this profile has already been computed using other 
methods than the ones shown here~\cite{1999PhRvD..60j5019B}. In this 
profile, there is a uniform field in the cylindrical region $\rho < \lambda$, and 
no field outside of the region:
\be 
B_z(\rho) = \frac{F}{\pi \lambda^2} \theta(\lambda-\rho).
\ee 
The profile function which gives this field is

\be 
f_\lambda(\rho)=\frac{\rho^2}{\lambda^2}\theta(\lambda-\rho) + \theta(\rho-\lambda)
\ee 
and the potential function is

\ba 
V_{m_l}(\rho)&=&\frac{m_l^2-1}{\rho^2} + \theta(\lambda-\rho)\frac{\mathcal{F}}{\lambda^2}
	\left(2(\sigma^3-m_l)+\frac{\mathcal{F}}{\lambda^2}\rho^2\right) \nonumber \\
	& &+ \theta(\rho-\lambda)\frac{\mathcal{F}}{\rho^2}
	\left(\mathcal{F}-2m_l\right) -\frac{2\mathcal{F}}{\lambda^2}\delta(\lambda-\rho).
\ea 

The solutions to equation (\ref{eqn:kchi}) for the inner region ($\rho < \lambda$) of 
this potential are expressed in terms of 
Whittaker $M$ and $W$ functions~\cite{slater1960confluent},

\be 
\label{eqn:Whit1}
[u_0(\rho)]_{\csubscripts}=\frac{W_{\frac{\lambda^2k^2}{4\mathcal{F}},\frac{m_l}{2}}
	\left(\frac{\mathcal{F}}{\lambda^2}\rho^2\right)}{\rho}
\ee 

\be 
\label{eqn:Whit2}
[u_\infty(\rho)]_{\csubscripts}=\frac{M_{\frac{\lambda^2k^2}{4\mathcal{F}},\frac{m_l}{2}}
	\left(\frac{\mathcal{F}}{\lambda^2}\rho^2\right)}{\rho},
\ee 
where 

\be
	\label{eqn:kchi}
	k^2=\chi^2-m^2-\frac{2\mathcal{F}}{\lambda^2}(\sigma^3-m_l).
\ee
The constant $W_0=\rho W(\rho)$ associated with these solutions is~\cite{slater1960confluent}

\be 
W_0^{-1}=-\frac{\lambda^2}{2\mathcal{F}}\frac{\Gamma\left[\frac{1}{2}\left(m_l+1-\frac{k^2\lambda^2}{2\mathcal{F}}\right)\right]}{\Gamma(1+m_l)}.
\ee 

The solutions for the exterior region ($\rho > \lambda$) are the Bessel functions,
\be 
\label{eqn:u0ext}
[u_0(\rho)]_{\csubscripts}=J_n(\sqrt{\chi^2-m^2}\rho)
\ee 
and
\be 
\label{eqn:uinfext}
[u_\infty(\rho)]_{\csubscripts}=Y_n(\sqrt{\chi^2-m^2}\rho),
\ee 
where $n=m_l-\mathcal{F}$.  We note that for large $m_l$ (i.e. $m_l\gg \mathcal{F}$), $n\approx m_l$ and these solutions coincide with the field-free case.

\subsection{Exterior Integral}
\label{sec:extint}

The homogeneous solutions for the exterior ($\rho > \lambda$) region are 
given by equations (\ref{eqn:u0ext}) and (\ref{eqn:uinfext}):

\ba 
	\Gamma^{(1)}_{\rm ext}& =& -\frac{\hbar\pi^2}{2} \sum_{\rm \sigma^3= \{\pm1\}}
	\sum_{m_l=-\infty}^{\infty}  
	\int_{0}^\infty \chi^3 d\chi
	\int_0^{\lambda} d\rho \rho^2 \biggr(\nonumber \\
	& &~~~~~~~J_{m_l-\mathcal{F}}(\sqrt{\chi^2 -m^2}\rho)
	Y_{m_l-\mathcal{F}}(\sqrt{\chi^2 -m^2}\rho)\nonumber \\ 
	& &~~~~~~~-J_{m_l}(\sqrt{\chi^2 -m^2}\rho)
	Y_{m_l}(\sqrt{\chi^2 -m^2}\rho) \biggr).
\ea
This integral is somewhat problematic because it is apparently divergent for large 
values of $\sqrt{\chi^2-m^2}\rho$. The asymptotic expansions of the Bessel functions 
for large argument are
\be
	J_n(x) \approx \sqrt{\frac{2}{\pi x}} \cos{\left(x-\frac{n\pi}{2}-\frac{\pi}{4}\right)}
\ee
and
\be
	Y_n(x) \approx \sqrt{\frac{2}{\pi x}} \sin{\left(x-\frac{n\pi}{2}-\frac{\pi}{4}\right)}.
\ee
So, in this limit, the integrand for an individual term of the $m_l$ sum results in a 
divergent integral:
\ba
	\lefteqn{J_{m_l-\mathcal{F}}(\sqrt{\chi^2 -m^2}\rho)
	Y_{m_l-\mathcal{F}}(\sqrt{\chi^2 -m^2}\rho)}\nonumber \\
	& &-
	J_{m_l}(\sqrt{\chi^2 -m^2}\rho)
	Y_{m_l}(\sqrt{\chi^2 -m^2}\rho) \nonumber \\
	& \approx&\frac{2}{\pi \sqrt{\chi^2-m^2}\rho}\sin{\left(\frac{\pi}{2}\mathcal{F}\right)}
	\cos{\left(\frac{\pi}{2}(\mathcal{F}-2m_l+1)+2\sqrt{\chi^2-m^2}\rho\right)}.
\ea
This seems to be an ultraviolet divergence, but it cannot be easily renormalized 
by a field-strength renormalization because the dependence on the field 
strength, $\mathcal{F}$, is not quadratic like the classical term. Nevertheless, 
since \ac{QED} is a renormalizable theory, the effective action in the exterior 
region must either be finite, or proportional to the classical action. To make 
sense of this divergence, we need to consider the entire sum over $m_l$, and 
not simply the individual terms in isolation.

We note that the exterior integral must be zero when 
$\mathcal{F}$ is an integer. In this case, the Bessel function order numbers 
of the first term are shifted 
by an integer value relative to the second term. But, since the sum is over all integer 
values of $m_l$, there will still be cancellation in pairs. In other words, for any
finite integer, $j$,
\be
	\sum_{i=-\infty}^{\infty}J_i(x)Y_i(x) = \sum_{i=-\infty}^{\infty}J_{i-j}(x)Y_{i-j}(x).
\ee
Physically, integer values of $\mathcal{F}$ correspond to the disappearance of the 
Aharonov-Bohm effect when the Berry's phase of a closed path is a multiple of $2\pi$. 
Since our primary motivation is in flux tubes in superconducting materials where 
$\mathcal{F}$ is naturally quantized to integer values, 
we will focus on integer values of $\mathcal{F}$ for now, and set this 
integral to zero. 

Edit: While the index-shifting argument correctly identifies the cancellation of the topological vacuum shift for integer flux, it is incomplete for a flux tube of finite radius. The magnetic core at $\rho < \lambda$ acts as a scattering center that modifies the exterior Green's function through the matching conditions at the boundary. As a result, the exterior density is non-zero even for integer flux, representing the physical diffusion of vacuum polarization into the field-free region. To correctly account for this, one must subtract a background Green's function that matches the local vector potential $A_\phi$ exactly at each point $\rho > \lambda$, ensuring that only the scattering contribution from the magnetic field—and not the spurious coordinate-dependent shift—is integrated.

The bulk of the effective action density is expected to be 
in the interior region where the field is, and 
where all of the classical action density is. However, nonlocal quantum effects can
cause diffusion of the effective action density to regions where 
no classical field exists~\cite{Gies:2001zp}. So, there may be interesting 
physical content within this integral.
Nevertheless, evaluating this integral for non-integer values of $\mathcal{F}$ is 
left for future research.

\subsection{Interior Integral}

Here, we consider the integral for the interior ($\rho < \lambda$) region.  
The Whittaker functions in equations (\ref{eqn:Whit1}) and (\ref{eqn:Whit2}) obey the differential
equation

\be 
\frac{d^2 W(\rho)}{d\rho^2} - \frac{1}{\rho}\frac{dW(\rho)}{d\rho} +\left[
	k^2+\frac{(1-m_l^2)}{\rho^2}-\rho^2\left(\frac{\mathcal{F}}{\lambda^2}\right)^2\right]W(\rho)=0.
\ee 
In exploring the asymptotics of the solutions, it is interesting to write this equation in 
canonical form

\be 
\label{eqn:canonical}
y''(\rho) +\left[k^2+\frac{1/4-m_l^2}{\rho^2} -\rho^2\left(\frac{\mathcal{F}}{\lambda^2}\right)^2 \right]y(\rho)=0,
\ee 
with $W(\rho)=\sqrt{\rho}y(\rho)$.

We may neglect the contribution from negative values of $m_l$ since for negative integer values of $n$~\cite{slater1960confluent},

\be 
	\frac{W_{k,\frac{n}{2} }(x) M_{k,\frac{n}{2}}(x)}{\Gamma(1+n)} = 0.
\ee 
For large values of $m_l$ ($m_l \gg \mathcal{F}$), the Whittaker functions correspond to Bessel 
functions and there is cancellation between them and the background terms.  So, we may compute the 
interior integral to arbitrary precision by summing a finite number of positive-$m_l$ terms.

\subsubsection{Field-Strength Renormalization}

There is an ultraviolet divergence for large values of $\chi$ which may be subtracted using a 
field-strength renormalization.  To see this, we consider the differential equation, (\ref{eqn:canonical}), 
in the limit of asymptotically large $k$ using the WKB approximation.  Our asymptotic $u(\rho)$ functions 
in this approximation are

\be 
\label{eqn:WKBu0}
[u_0(\rho)]_{k,m_l,\sigma^3} \approx \frac{1}{\sqrt{\rho}}\cos\left\{k\rho-\frac{1}{2k\rho}\left[(1/4-m_l^2) +
	\frac{\rho^4}{3}\left(\frac{\mathcal{F}}{\lambda^2}\right)^2\right]\right\}
\ee 
and
\be 
\label{eqn:WKBui}
[u_\infty(\rho)]_{k,m_l,\sigma^3} \approx \frac{1}{\sqrt{\rho}}\sin\left\{k\rho-\frac{1}{2k\rho}\left[(1/4-m_l^2) +
	\frac{\rho^4}{3}\left(\frac{\mathcal{F}}{\lambda^2}\right)^2\right]\right\}.
\ee 

The relevant constant for these solutions is $W_0=\rho W[u_0(\rho),u_\infty(\rho)]$, where

\be 
\label{eqn:WKBW0}
W_0\approx k+\frac{(1/4-m_l^2)}{2k\rho^2}-\frac{\rho^2}{2k}
	\left(\frac{\mathcal{F}}{\lambda^2}\right)^2.
\ee 
Taking equations (\ref{eqn:WKBu0}), (\ref{eqn:WKBui}), and 
(\ref{eqn:WKBW0}), and expanding in powers of $\frac{1}{k\rho}$ gives
\ba 
\lefteqn{\left[\frac{u_0(\rho)u_\infty(\rho)}{W_0}\right]_{k,m_l,\sigma^3} \approx
	\frac{\sin{(k\rho)}\cos{(k\rho)}}{k\rho}}\nonumber \\
	& &-\frac{1}{2(k\rho)^2}
	\left[(1/4-m_l^2)+\frac{1}{3}\left(\frac{\rho^2\mathcal{F}}{\lambda^2}\right)^2\right]\cos{(2k\rho)} \nonumber \\
	& & -\frac{(1/4-m_l^2)}{2(k\rho)^3}\sin{(k\rho)}\cos{(k\rho)} \nonumber \\
	& &+\frac{1}{2(k\rho)^3}\left(\frac{\rho^2\mathcal{F}}{\lambda^2}\right)^2\sin{(k\rho)}\cos{(k\rho)} \nonumber \\
	& & +\frac{1}{4(k\rho)^4}\left[(1/4-m_l^2)+\frac{1}{3}\left(\frac{\rho^2\mathcal{F}}{\lambda^2}\right)^2\right] \nonumber \\
	& &\times \left[(1/4-m_l^2)-\left(\frac{\rho^2\mathcal{F}} {\lambda^2}\right)^2\right]\cos{(2k\rho)} \\
	\label{eqn:infcn}
	\lefteqn{\approx \left[\frac{u_0(\rho)u_\infty(\rho)}{W_0}\right]_{k,m_l,\sigma^3} 
	\biggr|_{\mathcal{F}=0} }\nonumber\\
	 & &-\frac{1}{6(k\rho)^2}\left(\frac{\rho^2\mathcal{F}}
	 {\lambda^2}\right)^2\cos{(2k\rho)}\nonumber \\
	 & & +\frac{1}{2(k\rho)^3}
	\left(\frac{\rho^2\mathcal{F}}
	 {\lambda^2}\right)^2\sin{(k\rho)}\cos{(k\rho)} \nonumber \\
	 & & -\frac{1}{6(k\rho)^4}\left(\frac{\rho^2\mathcal{F}}
	 {\lambda^2}\right)^2\biggr[(1/4-m_l^2)  -	\left(\frac{\rho^2\mathcal{F}}
	 {\lambda^2}\right)^2\biggr]\cos{(2k\rho)}.
\ea 


The subtraction of the background field corresponds to cancellation of the first term in 
equation (\ref{eqn:infcn}).  The remaining terms represent ultraviolet divergences proportional 
to $\left(\frac{\mathcal{F}}{\lambda^2}\right)^2$.  These terms may be subtracted out from the integral
 by a field-strength renormalization.

The renormalized contribution to the effective action from the interior region is, then

%\ba 
%	\label{eqn:gfee}
%	\Gamma^{(1)}_{\rm int}& =& -\hbar\pi \sum_{\rm \sigma^3= \{\pm1\}}\sum_{m_l=0}^{m_l\gg\mathcal{F}}  
%	\int_{0}^\infty \chi^3 d\chi
%	\int_0^{\lambda} d\rho \rho^2 \biggr(
%	\frac{\lambda^2}{2\mathcal{F}}\nonumber \\
%	& &\frac{\Gamma\left(\frac{1}{2}\left(m_l+1-\frac{k^2\lambda^2}{2\mathcal{F}}\right)\right)}
%		{m_l!} 
%	W_{\frac{\lambda^2k^2}{4\mathcal{F}},\frac{m_l}{2}}
%	\left(\frac{\mathcal{F}}{\lambda^2}\rho^2\right) 
%	M_{\frac{\lambda^2k^2}{4\mathcal{F}},\frac{m_l}{2}}
%	\left(\frac{\mathcal{F}}{\lambda^2}\rho^2\right)\nonumber \\ 
%	& &~~~~~~~-\frac{\pi}{2}J_{m_l}(\sqrt{\chi^2 -m^2}\rho)
%	Y_{m_l}(\sqrt{\chi^2 -m^2}\rho)\nonumber \\
%	& & +\left(\frac{\mathcal{F}}{\lambda^2}\right)^2
%	\biggr[ \frac{\rho^3}{k^2} \sin(\Theta_{m_l,k}(\rho))
%	\cos(\Theta_{m_l,k}(\rho))+\nonumber \\
%	& &\frac{\rho^2}{6k^3}\cos(2\Theta_{m_l,k}(\rho))\biggr]\biggr).
%\ea 

\ba 
	\label{eqn:gfee}
	\Gamma^{(1)}_{\rm int}& =& -\hbar\pi \sum_{\rm \sigma^3= \{\pm1\}}\sum_{m_l=0}^{m_l\gg\mathcal{F}}  
	\int_{0}^\infty \chi^3 d\chi
	\int_0^{\lambda} d\rho \rho^2 \biggr(
	\frac{\lambda^2}{2\mathcal{F}}\nonumber \\
	& &\frac{\Gamma\left(\frac{1}{2}\left(m_l+1-\frac{k^2\lambda^2}{2\mathcal{F}}\right)\right)}
		{m_l!} 
	W_{\frac{\lambda^2k^2}{4\mathcal{F}},\frac{m_l}{2}}
	\left(\frac{\mathcal{F}}{\lambda^2}\rho^2\right) 
	M_{\frac{\lambda^2k^2}{4\mathcal{F}},\frac{m_l}{2}}
	\left(\frac{\mathcal{F}}{\lambda^2}\rho^2\right)\nonumber \\ 
	& &~~~~~~~-\frac{\pi}{2}J_{m_l}(\sqrt{\chi^2 -m^2}\rho)
	Y_{m_l}(\sqrt{\chi^2 -m^2}\rho)\nonumber \\
	& & +\left(\frac{\mathcal{F}}{\lambda^2}\right)^2
	\biggr[ \frac{\rho^3}{2k^2} \sin(\Theta_{m_l,k}(\rho))
	+\frac{\rho^2}{6k^3}\cos(\Theta_{m_l,k}(\rho))\biggr]\biggr),
\ea 
where $\Theta_{m_l,k}(\rho)\equiv 2k\rho -\frac{(1/4-m_l^2)}{k\rho}$ and $k$ is given 
in terms of the relevant integration and summation parameters by 
equation (\ref{eqn:kchi}).



\section{Discussion}
Equation (\ref{eqn:gfee}) is a renormalized expression of the 1-loop correction term to the effective 
action for a step-function flux tube in terms of finite integrals and sums over finitely many terms. 
So, it is straightforward to implement a numerical algorithm for computing the effective action using 
this expression. However,  the integrand is polluted by numerous poles
which make the computations time consuming. For this reason, it is more practical to approach the problem 
using other numerical methods such as those presented in chapter \ref{ch:WLNumerics}. 
Moreover, the effective action for the step function flux tube profile 
has been computed before, arguably in a simpler way~\cite{1999PhRvD..60j5019B}. 
Establishing if the method discussed here produces the same results as 
previous methods is an important next step along this research path.
Nevertheless, the Green's function 
method described in this chapter has proved useful for 
establishing an expression for non-homogeneous 
configurations with stationary effective action, equation 
(\ref{eqn:stationaryEffAct}). Additionally, it may 
prove useful in future research since it is a technique which 
can be applied to arbitrary flux tube 
profiles, and doesn't have the precision limitations of Monte Carlo 
techniques. One could also imagine that 
a clever choice of potential might lead to considerable analytic 
progress. For example, some 
potential function 
may result in functions $u_0(\rho)$ and $u_\infty(\rho)$ for 
which the integrals can be evaluated.

One obstacle with the approach outlined in this chapter is in performing 
the field-strength renormalization in general. In section \ref{sec:extint}, we 
observed that the renormalization proceedure could not be performed on individual 
terms alone, but rather needed to account for the entire angular momentum summation.
This poses certain technical problems that were not addressed here. However, previous 
research applying a WKB phase shift method to fermion determinants in non-homogeneous 
background fields has produced a renormalization scheme which may be generalized to 
the background fields discussed in this chapter~\cite{2004PhLB..600..302D}. 
There are therefore promising prospects in addressing the field-strength renormalization 
in this method by proceeding analogously with the WKB phase shift method.

Perhaps the most interesting result from this section is that 
the Green's function method may provide 
a way to numerically compute the
quantum corrected equations of motion for external 
electromagnetic fields. Making further progress on studying 
equation (\ref{eqn:stationaryEffAct}) may therefore be a
particularly interesting direction for future research.
